% Part: computability
% Chapter: recursive-functions
% Section: other-recursions

\documentclass[../../../include/open-logic-section]{subfiles}

\begin{document}

\olfileid{cmp}{rec}{ore}
\olsection{نور بازګښتونه}

د جوړو او لړيو په کارولو د بنسټيز بازګښت لا نابلده (او ګټورې) بڼې توجيه کولے شو۔ مثلاً ډېر وخت دوه تابعې يوځاے تعريفول ګټور وي، لکه په لاندې تعريف کښې:
\begin{align*}
h_0(\vec x, 0) & = f_0(\vec x) \\
h_1(\vec x, 0) & = f_1(\vec x) \\
h_0(\vec x, y+1) & = g_0(\vec x, y, h_0(\vec x, y), h_1(\vec x, y)) \\
h_1(\vec x, y+1) & = g_1(\vec x, y, h_0(\vec x, y), h_1(\vec x, y))
\end{align*}
دا د \emph{هممهاله بازګښت} يوه بېلګه ده۔ د تابعو د تعريف بله ګټوره لار دا ده چې د~$h(\vec x, y+1)$ ارزښت د ټولو ارزښتونو $h(\vec x, 0)$، \dots،~$h(\vec x, y)$ له مخې ورکړو، لکه په لاندې تعريف کښې:
\begin{align*}
h(\vec x, 0) & = f(\vec x) \\
h(\vec x, y+1) & = g(\vec x, y, \tuple{h(\vec x, 0), \dots, h(\vec x, y)}).
\end{align*}
لاندې چوکاټ دا مفکوره لا لنډه بيانوي:
\[
h(\vec x, y) = g(\vec x, y, \tuple{h(\vec x, 0), \dots, h(\vec x, y-1)})
\]
په دې پوهېدو سره چې کله~$y$، $0$ وي، د~$g$ وروستے آرګومېنټ يوازې تشه لړۍ وي۔ په دواړو بڼو کښې مفکوره دا ده چې د «راتلونکي ګام» د محاسبې پر مهال تابع~$h$ تر دې ځايه د محاسبه شوو ارزښتونو ټوله لړۍ کارولے شي۔ دې ته \emph{د ارزښتونو د بهير بازګښت} وايي۔ د يوې ځانګړې بېلګې دپاره لاندې ډول تعريف توجيه کولے شي:
\begin{align*}
h(\vec x, y) & = \begin{cases}
  g(\vec x, y, h(\vec x, k(\vec x, y))) & \text{که $k(\vec x, y) < y$ وي} \\
  f(\vec x) & \text{که نۀ}
\end{cases}
\end{align*}
په بله وينا، په~$y$ د~$h$ ارزښت د~$h$ د هر مخکيني ارزښت له مخې محاسبه کېدے شي، او هغه مخکينی ارزښت~$k$ ورکوي۔

\begin{prob}
  د پاتې شوني تابع~$r(x,y)$ د ارزښتونو د بهير په بازګښت تعريف کړئ۔ (که~$x$ او~$y$ طبيعي عددونه او~$y > 0$ وي، نو~$r(x,y)$ له~$y$ څخه هغه کوچنے عدد دے چې د کوم~$z$ دپاره $x = z\times y + r(x,y)$ وي۔ د کره‌والي دپاره وايو که~$y=0$ وي، نو~$r(x,0) = 0$۔)
\end{prob}

بايد په دې فکر وکړئ چې دا تابعې په عادي بنسټيز بازګښت څنګه ترلاسه کېږي۔ د بنسټيز بازګښت يوه وروستۍ بڼه په دې لحاظ لا انعطاف منونکې ده چې د لارې په اوږدو کښې د \emph{پاراميټرونو} (اړخيزو ارزښتونو) بدلولو اجازه ورکوي:
\begin{align*}
h(\vec x, 0) & = f(\vec x) \\
h(\vec x, y+1) & = g(\vec x, y, h(k(\vec x), y))
\end{align*}
دا هم په عادي بنسټيز بازګښت تقليدېدے شي۔ (دا کول ستونزمن دي۔ د يوې لارښوونې دپاره، محاسبه په لاس له پاے څخه بېرته پرانيزئ۔)

\end{document}
